% sections/01_models_literature.tex  (owner: lit)
% Literature survey and novelty check. References: bib/lit.bib.
% Literature state as found by web searches on 2026-10-13.
\section{Industrial wave models with self-similar blowup: a survey}
\label{lit:sec}

All models collected below are, after nondimensionalisation, of the form
\begin{equation}\label{lit:eq:generic}
  \varepsilon^{2}\partial_t^2 u+\gamma\,\partial_t u-\Delta u+m^{2}u
  =F(u)\quad\text{or}\quad F(u,\partial u),
\end{equation}
possibly with a fourth-order or quasilinear principal part. Near a
singularity at $(T,x_0)$ the time scale $T-t$ is short compared with the
relaxation time $\varepsilon^{2}/\gamma$, so after rescaling $t$ the
principal part is the wave operator and $\gamma\partial_t u$, $m^2u$ are
lower order (in the similarity variables of \cref{fw:sec} they carry
the factors $e^{-\tau}$ and $e^{-2\tau}$). The prototype singularity is the
spatially homogeneous (``ODE'') solution $u_T$ of $u''=F(u)$, which is
self-similar: $u_T=\kappa(T-t)^{-s}$ or $u_T=-2\log(T-t)+\mathrm{const}$.
We say \emph{blowup} if $|u|\to\infty$ and \emph{quenching} (in MEMS:
\emph{touchdown}) if $u$ stays bounded while $F(u)$ and $\partial_tu$ become
unbounded; in similarity variables the two cases differ only in the sign of
the scaling exponent $s$ ($s<0$ for quenching). For each model class we
record the equation and its industrial origin, known blowup/quenching
results, self-similar solutions and stability results; a summary is given
in Table~\ref{lit:tab:summary}.

\subsection{Electrostatic MEMS with inertia: touchdown as quenching}
\label{lit:mems}
In the small-aspect-ratio limit, the deflection $u\in(-1,0]$ of an
electrostatically actuated elastic membrane suspended above a rigid ground
plate (gap $1+u$) satisfies
\begin{equation}\label{lit:eq:mems}
  \varepsilon^{2}u_{tt}+u_t-\Delta u=-\frac{\lambda f(x)}{(1+u)^{2}}
  \quad\text{in }\Omega,\qquad u=0\ \text{on }\partial\Omega,
\end{equation}
where $\varepsilon^2$ is the ratio of inertial to damping forces,
$\lambda\propto V^2$ measures the applied voltage and $f$ the dielectric
profile; see the monographs
\cite{PeleskoBernstein2002,EspositoGhoussoubGuo2010,KavallarisSuzuki2018},
the survey \cite{LaurencotWalker2017}, and
\cite{FloresEtAl2007,Flores2014} for the hyperbolic model
\eqref{lit:eq:mems}. Variants include a nonlocal version, in which $\lambda$
is replaced by $\lambda/(1+\alpha\int_\Omega(1+u)^{-1}\,dx)^{2}$ (device in
series with a capacitor) \cite{KavallarisEtAl2011,KavallarisSuzuki2018}, a
fringing-field correction \cite{Miyasita2017}, and plate models with
$\Delta^2u$ \cite{LaurencotWalker2014,LaurencotWalker2017}. At nanometre
gaps, van der Waals and Casimir forces (scaling like $w^{-3}$ and $w^{-4}$
in the gap $w$) compete with the electrostatic force $w^{-2}$ and can cause
collapse even without applied voltage
\cite{SerryWalliserMaclay1998,BatraPorfiriSpinello2007}. With $w:=1+u$ and
in the inertia-dominated regime one is led to the model
\begin{equation}\label{lit:eq:memsq}
  w_{tt}-\Delta w=-\lambda\,w^{-q},\qquad q>1\quad
  (q=2 \text{ electrostatic},\ q=3 \text{ van der Waals},\ q=4\text{ Casimir}).
\end{equation}

\emph{ODE quenching.} Equation \eqref{lit:eq:memsq} has the self-similar
solution
\begin{equation}\label{lit:eq:memsode}
  w_T(t)=\kappa_q\,(T-t)^{\sigma},\qquad \sigma=\tfrac{2}{q+1},\qquad
  \kappa_q^{\,q+1}=\frac{\lambda(q+1)^{2}}{2(q-1)}
  \qquad(q=2:\ \sigma=\tfrac23,\ \kappa_2^3=\tfrac{9\lambda}{2}),
\end{equation}
for which $w_T\to0$ and $\partial_tw_T\to-\infty$ as $t\to T$. In the
notation of \cref{fw:sec}, $s=-\sigma$ and $V_0=F'(w_T)(T-t)^2=2q(q-1)/(q+1)^2$; the
roots of $a^2+a-V_0=0$ are $a_1=1-\sigma$ and $a_2=-(2-\sigma)$, so the
candidate eigenvalues are $1-\ell-2n$ and $-2(q-1)/(q+1)-\ell-2n$. The only
non-negative eigenvalues are $1$ (time translation) and $0$ with $\ell=1$
(Lorentz boosts), exactly as for focusing power nonlinearities.
(This is the formal count; it is verified in the model sections.)

\emph{Known results.} Local well-posedness, global existence for small
$\lambda$ and finite-time touchdown for large $\lambda$ are known for
\eqref{lit:eq:mems} and its variants
\cite{ChangLevine1981,FloresEtAl2007,Guo2010,KavallarisEtAl2011,Flores2014,KavallarisEtAl2015,Miyasita2017};
\cite{Flores2014} shows numerically that the dynamic pull-in threshold lies
below the static one and depends on $\varepsilon$, and
\cite{KavallarisEtAl2015} studies the quenching behaviour of
$u_{tt}=u_{xx}+\lambda(1-u)^{-2}$ in one dimension. For
\eqref{lit:eq:memsq} with radial data, Gimperlein, He and Lacey
\cite{GimperleinHeLacey2022} (see also \cite{He2021}) obtain
\emph{formally} a sequence of self-similar quenching solutions; formal
asymptotics and numerics indicate that radial solutions increasing in $|x|$
quench at the origin like \eqref{lit:eq:memsode}, that the trivial profile
is stable up to the time-translation mode, and that the non-trivial
profiles are unstable. For the parabolic MEMS problem, quenching profiles
are constructed rigorously in self-similar variables, e.g.\
\cite{LiaoNguyen2025}. We found \emph{no} rigorous stability result for the
hyperbolic touchdown \eqref{lit:eq:memsode} (see
Section~\ref{lit:novelty}).

\subsection{Hyperbolic heat conduction and thermal explosion}
\label{lit:heat}
The Maxwell--Cattaneo--Vernotte law $\tau\boldsymbol q_t+\boldsymbol
q=-k\nabla\theta$ \cite{Cattaneo1958,Vernotte1958} (reviews:
\cite{JosephPreziosi1989,Straughan2011}) replaces Fourier's law when the
time scale is comparable with the thermal relaxation time $\tau$ (pulsed
laser heating, fast reactions). Combined with the energy balance $\rho
c\,\theta_t+\nabla\!\cdot\boldsymbol q=Q(\theta)$ with an Arrhenius or
Frank-Kamenetskii source $Q(\theta)=\delta e^{\theta}$ (thermal explosion,
combustion \cite{BebernesEberly1989}), elimination of $\boldsymbol q$ gives
\begin{equation}\label{lit:eq:mcv}
  \tau\rho c\,\theta_{tt}+\rho c\,\theta_t-k\Delta\theta
  =Q(\theta)+\tau Q'(\theta)\,\theta_t .
\end{equation}
The ``hyperbolic Gelfand problem'' usually studied is the semilinear model
\begin{equation}\label{lit:eq:gelfand}
  u_{tt}+\gamma u_t-\Delta u=\delta e^{u},\qquad
  u_T(t)=\log\tfrac{2}{\delta}-2\log(T-t),
\end{equation}
i.e.\ \eqref{lit:eq:mcv} without the last term. A short computation shows
that this omission matters: for the full system the flux of a spatially
homogeneous state is divergence free, so homogeneous solutions obey the
first-order ODE $\rho c\,\theta_t=Q(\theta)$ and blow up like
$\theta\approx-\log(T-t)$ (parabolic Frank-Kamenetskii type), and near such
a state $\tau Q'(\theta)\theta_t$ is of the same order as
$\tau\rho c\,\theta_{tt}$. Thus the wave-type ODE blowup in
\eqref{lit:eq:gelfand} is a feature of the semilinear model (and of
Liouville-type equations), not of the Cattaneo system itself; we did not
find the singularity formation of \eqref{lit:eq:mcv} treated in the
literature.

For \eqref{lit:eq:gelfand}: in $1+1$ dimensions $u_{tt}-u_{xx}=e^u$ is the
classical Liouville equation, whose general solution is explicit.
Kichenassamy \cite{Kichenassamy1996} constructs, by Fuchsian reduction and a
Nash--Moser argument and in any dimension, solutions blowing up on a
prescribed spacelike hypersurface with logarithmic behaviour, and according
to the abstract treats Cauchy data close to constants (we have only seen
the abstract). In one dimension, Godin \cite{Godin2000} shows for mixed
problems that solutions are either global or blow up on a spacelike curve,
and Azaiez--Masmoudi--Zaag \cite{AzaiezMasmoudiZaag2016} determine the
blowup rate at non-characteristic points for low-regularity data.

\subsection{Softening foundations and focusing Klein--Gordon equations}
\label{lit:kg}
A string or beam on a softening (Duffing-type) elastic foundation, after
scaling, gives
\begin{equation}\label{lit:eq:kg}
  u_{tt}-u_{xx}+u=|u|^{p-1}u\quad(p=3)\qquad\text{or}\qquad
  u_{tt}+u_{xxxx}+u=u^{3},
\end{equation}
the focusing Klein--Gordon ($\phi^4$) equation in the string case. Since
the cubic law is a truncated Taylor expansion, blowup signals a loss of
stability (snap-through, collapse) rather than physically infinite
displacement. Fourth-order suspension-bridge models show finite-time blowup
with wide oscillations, which is not self-similar
\cite{GazzolaPavani2013,Gazzola2015}. For \eqref{lit:eq:kg} the ODE
profile is $\kappa(T-t)^{-s}$, $s=2/(p-1)$, $\kappa^{p-1}=s(s+1)$
($\sqrt2/(T-t)$ for $p=3$), and the mass term is an $e^{-2\tau}$
perturbation. Blowup rates for focusing NLKG are due to Killip, Stovall and
Visan \cite{KillipStovallVisan2014} and Hamza--Zaag \cite{HamzaZaag2013}
(one dimension or radial). Donninger and Wallauch \cite{DonningerWallauch2021}
show for a class of strongly perturbed supercritical equations in $d=3$
that the ODE profile of the unperturbed equation still describes stable
blowup. The damped analogue $u_{tt}+\gamma u_t-\Delta u=|u|^{p-1}u$ is
covered by Lyapunov-functional methods in similarity variables, which give
the blowup rate for lower-order perturbations
\cite{HamzaZaag2010,HamzaZaag2012}; in one dimension, solutions converging
exponentially to the ODE profile are constructed in \cite{HamzaSaidi2021}.

\subsection{Nonlinear acoustics: derivative-quadratic wave equations}
\label{lit:acoustics}
High-intensity focused ultrasound (HIFU) and lithotripsy are modelled by the Westervelt and Kuznetsov equations
\cite{Kaltenbacher2015}
\begin{equation}\label{lit:eq:west}
  p_{tt}-c^2\Delta p-b\Delta p_t=k\,(p^2)_{tt},\qquad
  \psi_{tt}-c^2\Delta\psi-b\Delta\psi_t
  =\partial_t\big(|\nabla\psi|^2+\tfrac{B}{2A}c^{-2}\psi_t^2\big).
\end{equation}
With viscosity $b>0$ there is global well-posedness for small data
\cite{KaltenbacherLasiecka2009}. In the inviscid case \eqref{lit:eq:west}
is a quasilinear wave equation with local well-posedness and an expected
gradient blowup (shocks) \cite{DekkersRozanovaPierrat2017}. Two ODE-type
mechanisms are relevant. (i) \emph{Degeneracy}: for $b=0$ the Westervelt
equation reads $\partial_t^2(p-kp^2)=c^2\Delta p$; spatially homogeneous
solutions have $p-kp^2=A+Bt$, so for $B>0$ the branch with $p<1/(2k)$ satisfies
$1-2kp=\sqrt{4kB(T-t)}$, which vanishes
at a finite time $T$ with $p_t\sim(T-t)^{-1/2}$: $p$ stays bounded while
the equation degenerates, a quasilinear analogue of quenching (our
observation; we found no study of it). (ii) \emph{Riccati-type
nonlinearities}: the semilinear model $u_{tt}-\Delta u=(\partial_tu)^2$
has the ODE blowup $u_T=-\log(T-t)$. Speck \cite{Speck2020} proves stable
ODE-type blowup in $d=3$ for a class of quasilinear equations with such
derivative-quadratic terms. In one dimension, Gough \cite{Gough2025} and
Liu--Raees \cite{LiuRaees2025} show that no smooth exact self-similar
profiles exist, construct generalised self-similar solutions with
logarithmic growth, and prove asymptotic stability of the relevant branches
(including the ODE branch). The spatial-derivative analogue is treated in
\cite{GhoulLiuMasmoudi2025}. Note that $u_{tt}-\Delta u=u_t^2-|\nabla u|^2$
is linearised by $v=e^{-u}$ ($\Box v=0$), so its blowup is the touchdown
$v\to0$ of a free wave.

\subsection{Nematic liquid crystals, \texorpdfstring{$\sigma$}{sigma}-models and wave maps}
\label{lit:lc}
If the inertia of the director field $\boldsymbol n$ is retained and the
elastic anisotropy and the coupling to flow are neglected, a nematic
director $\boldsymbol n:\mathbb{R}^{1+3}\to S^2$ formally obeys the wave
map equation $\boldsymbol n_{tt}-\Delta\boldsymbol n=(|\nabla\boldsymbol
n|^2-|\boldsymbol n_t|^2)\boldsymbol n$. With anisotropic (splay/bend)
elasticity a planar director angle $\psi$ satisfies the variational wave
equation $\psi_{tt}-c(\psi)(c(\psi)\psi_x)_x=0$, for which Glassey, Hunter
and Zheng proved finite-time singularity formation (cusps; gradient
blowup, not ODE-type) and nonlinear instability of constant states
\cite{GlasseyHunterZheng1996}. Self-similar blowup for co-rotational wave
maps $\mathbb{R}^{1+3}\to S^3$ goes back to Shatah \cite{Shatah1988} and
Turok--Spergel \cite{TurokSpergel1990}, with excited profiles in
\cite{Bizon2000} and explicit profiles $\mathbb{R}^{1+d}\to S^d$, $d\ge4$, in
\cite{BiernatBizon2015}. These are the founding examples of the Donninger
programme (Section~\ref{lit:donninger}). The stability theory is developed
for targets $S^d$ with $d$ equal to the space dimension; we are not aware of
results specific to the $S^2$-valued director in three dimensions.

\subsection{Relativistic strings, membranes and Born--Infeld}
\label{lit:membranes}
Timelike extremal (zero mean curvature) hypersurfaces in Minkowski space
model relativistic strings and membranes and, in $1+1$ dimensions, give
the Born--Infeld equation
$(1+u_x^2)u_{tt}-2u_tu_xu_{tx}-(1-u_t^2)u_{xx}=0$. The origin is physical
rather than industrial. Eggers and Hoppe describe singularity formation by
self-similar solutions \cite{EggersHoppe2009}, extended to membranes in
\cite{EggersEtAl2015}. Yan constructs explicit self-similar solutions and
states their stability inside a strict sub-cone of the backward light cone
\cite{Yan2019a,Yan2019b}; by \cref{geom:prop:cone}(c), the explicit
solutions $\pm(T-t)\sqrt{1-|x|^2/(T-t)^2}$ of \cite{Yan2019b} are null
cones, on which the equation degenerates. These equations are quasilinear and lie
outside the semilinear framework below.

\subsection{Further models}
\label{lit:other}
MEMS plates, $u_{tt}+\Delta^2u=-\lambda(1+u)^{-2}$
\cite{LaurencotWalker2014,LaurencotWalker2017}, have no finite speed of
propagation, and their natural scaling is $x\sim(T-t)^{1/2}$, so the
light-cone method does not apply directly. Coupled MEMS models, in which
the electric potential solves an elliptic problem in the moving gap, are
analysed in \cite{He2021,LaurencotWalker2017}. Cattaneo-type (hyperbolic)
chemotaxis models are a further candidate. % TODO verify references (hyperbolic Keller--Segel) -- not cited
We have not reviewed their singularities.

\begin{table}[htbp]
\centering
\scriptsize
\setlength{\tabcolsep}{2pt}
\renewcommand{\arraystretch}{1.15}
\begin{tabular}{@{}p{0.11\linewidth}p{0.22\linewidth}p{0.12\linewidth}p{0.17\linewidth}p{0.14\linewidth}p{0.18\linewidth}@{}}
\toprule
\raggedright Model & \raggedright Equation & \raggedright Application & \raggedright Self-similar profile & \raggedright Symmetry modes & \raggedright Stability status \tabularnewline
\midrule
\raggedright Hyperbolic MEMS & \raggedright $w_{tt}+\gamma w_t-\Delta w=-\lambda w^{-q}$, $q=2,3,4$ & \raggedright MEMS/NEMS touchdown & \raggedright $\kappa_q(T-t)^{2/(q+1)}$ (quenching); formal non-trivial ones \cite{GimperleinHeLacey2022} & \raggedright $\lambda=1$ (time); $\lambda=0$, $\ell=1$ (boosts) & \raggedright Open in the literature (formal and numerical evidence only); target of this note \tabularnewline
\raggedright Nonlocal MEMS & \raggedright as above, $\lambda\mapsto\lambda(1+\alpha\!\int\! w^{-1})^{-2}$ & \raggedright MEMS with series capacitor & \raggedright as above (formally) & \raggedright as above & \raggedright Open \tabularnewline
\raggedright MEMS plate & \raggedright $w_{tt}+\Delta^2w=-\lambda w^{-2}$ & \raggedright MEMS plates & \raggedright parabolic scaling & \raggedright n/a & \raggedright Open; light-cone method not applicable \tabularnewline
\raggedright Hyperbolic Gelfand & \raggedright $u_{tt}+\gamma u_t-\Delta u=\delta e^u$ & \raggedright thermal explosion (semilinear model) & \raggedright $\log\frac2\delta-2\log(T-t)$ & \raggedright $\lambda=1$; $\lambda=0$, $\ell=1$ & \raggedright Smooth data: \cite{Kichenassamy1996}; 1D: \cite{Godin2000,AzaiezMasmoudiZaag2016}; $H^k$ light-cone stability open \tabularnewline
\raggedright Cattaneo system & \raggedright \eqref{lit:eq:mcv} & \raggedright laser heating, combustion & \raggedright $-\log(T-t)$ (parabolic type) & \raggedright $\lambda=1$; $\lambda=0$, $\ell=1$ (tilts, \cref{therm:sec:full}) & \raggedright Open (not found) \tabularnewline
\raggedright Focusing KG & \raggedright $u_{tt}-\Delta u+m^2u=|u|^{p-1}u$ & \raggedright softening foundations, $\phi^4$ & \raggedright $\kappa(T-t)^{-2/(p-1)}$ & \raggedright $\lambda=1$; $\lambda=0$, $\ell=1$ & \raggedright Rates \cite{KillipStovallVisan2014,HamzaZaag2013}; $m=0$: \cite{Ostermann2024}; $m>0$: expected \cite{DonningerWallauch2021} \tabularnewline
\raggedright Damped power & \raggedright $u_{tt}+\gamma u_t-\Delta u=|u|^{p-1}u$ & \raggedright damped structures & \raggedright as above & \raggedright as above & \raggedright Rates \cite{HamzaZaag2010,HamzaZaag2012}; 1D \cite{HamzaSaidi2021}; $H^k$, $d\ge2$ open \tabularnewline
\raggedright Riccati derivative & \raggedright $u_{tt}-\Delta u=(\partial_tu)^2$ & \raggedright nonlinear acoustics (model) & \raggedright $-\log(T-t)$; generalised SS in 1D & \raggedright $\lambda=1$ (time); $\lambda=0$ (shift $u\mapsto u+c$; tilts, $\ell=1$); no boosts & \raggedright 1D \cite{Gough2025,LiuRaees2025}; 3D (other method) \cite{Speck2020} \tabularnewline
\raggedright Westervelt, $b=0$ & \raggedright $\partial_t^2(p-kp^2)=c^2\Delta p$ & \raggedright HIFU, lithotripsy & \raggedright $1-2kp\propto(T-t)^{1/2}$ & \raggedright formal: $1$; $\frac14$ (tilt); $0$ (\cref{acou:sub-west}) & \raggedright Open (quasilinear) \tabularnewline
\raggedright Wave maps & \raggedright $\Box\boldsymbol n=-(|\partial_t\boldsymbol n|^2-|\nabla\boldsymbol n|^2)\boldsymbol n$ & \raggedright $\sigma$-models, nematics & \raggedright Shatah--Turok--Spergel, Biernat--Bizo\'n & \raggedright $\lambda=1$ & \raggedright Known (Section~\ref{lit:donninger}) \tabularnewline
\raggedright Variational wave & \raggedright $\psi_{tt}-c(\psi)(c(\psi)\psi_x)_x=0$ & \raggedright nematic LC & \raggedright cusp singularities & \raggedright n/a & \raggedright Instability of constants \cite{GlasseyHunterZheng1996} \tabularnewline
\raggedright Timelike minimal surfaces & \raggedright Born--Infeld & \raggedright strings, membranes & \raggedright \cite{EggersHoppe2009,EggersEtAl2015} & \raggedright n/a & \raggedright Partial (sub-cone, null-cone solutions) \cite{Yan2019b} \tabularnewline
\bottomrule
\end{tabular}
\caption{Wave-type models with self-similar singularities. ``Symmetry modes''
are the non-negative eigenvalues of the formal linearisation around the
ODE profile (\cref{fw:sec}); for co-rotational wave maps only time
translation survives the symmetry reduction. Lower-order terms ($\gamma$,
$m$) break Lorentz invariance but decay like $e^{-\tau}$, $e^{-2\tau}$.}
\label{lit:tab:summary}
\end{table}

\subsection{The Donninger et al.\ method: literature}
\label{lit:donninger}
The method writes the equation as a first-order system in similarity
coordinates on the backward light cone, so that the self-similar profile
becomes a static solution. The linearised operator is treated as a
compact perturbation of a generator $\mathbf L_0$ with explicit growth
bound, obtained by Lumer--Phillips in $H^k\times H^{k-1}$. One proves mode
stability, removes the symmetry-induced unstable eigenvalues by Riesz
projections, modulates the blowup time (and the boost parameters), and
closes a fixed-point argument. A survey is \cite{Donninger2024}. All
entries below were checked against arXiv and journal or repository records.
\begin{itemize}[leftmargin=1.2em,itemsep=1pt]
\item \emph{Wave maps and Yang--Mills, conditional and mode stability.}
\cite{Donninger2011}: nonlinear stability of the Shatah--Turok--Spergel
co-rotational wave map $\mathbb{R}^{1+3}\to S^3$, conditional on mode
stability. \cite{DonningerSchoerkhuberAichelburg2012}: the linearised
problem; it excludes unstable eigenvalues with large imaginary part or real
part $>\frac12$. \cite{CostinDonningerXia2016}: rigorous mode stability of
this wave map, completing the proof (note: wave maps, not Yang--Mills).
\cite{Donninger2014}: conditional stability for $(1+5)$-dimensional
equivariant Yang--Mills. \cite{CostinEtAl2016}: mode stability for that
Yang--Mills profile, via explicit approximate (``quasi-'') solutions; this
completes the stability proof. \cite{Glogic2022}: stable blowup for
Yang--Mills in odd $d\ge5$.
\cite{CostinDonningerGlogic2017}: mode stability of the Bizo\'n--Biernat
wave maps for $d\ge4$. \cite{DonningerGlogic2019}: existence and stability
of self-similar blowup for wave maps into negatively curved targets, $d\ge8$.
\item \emph{Power nonlinearities, ODE blowup.} \cite{DonningerSchoerkhuber2012}:
$1<p\le3$, $d=3$, radial, energy topology.
\cite{DonningerSchoerkhuber2014}: $p>3$, $d=3$, radial, stronger topology.
\cite{DonningerSchoerkhuber2016}: $d=3$ without symmetry, $H^2\times H^1$. \cite{DonningerSchoerkhuber2017}: odd $d\ge5$,
radial. \cite{Donninger2017}: energy-critical $p=5$, $d=3$, radial; Strichartz
estimates in similarity coordinates. \cite{DonningerRao2018}: the energy-critical
case at optimal regularity (preprint). \cite{Wallauch2022}: energy-critical,
$3\le d\le6$ (preprint). \cite{ChatzikaleasDonninger2019}: cubic, odd
$d\ge5$, without symmetry. \cite{Ostermann2024}: all $d$, all $p>1$, without
symmetry. \cite{DonningerWallauch2021}: strongly perturbed equations, $d=3$.
\cite{Liu2025}: superconformal ODE family in all dimensions (preprint).
\item \emph{Non-trivial self-similar profiles, finite codimension.}
\cite{GlogicSchoerkhuber2021}: explicit profile of the cubic equation for
$d\ge5$; codimension-one stability for $d=7$.
\cite{CsoboGlogicSchoerkhuber2024}: quadratic equation, explicit profile for
$d\ge7$, codimension-one stability for $d=9$, non-radial ODE stability for
$d=7,9$.
\item \emph{Beyond the light cone.} \cite{BiernatDonningerSchoerkhuber2021}:
hyperboloidal similarity coordinates; global-in-space stability of the
$S^3$ wave map. \cite{ChenEtAl2024}: quadratic equation, $d=7$, conditional
codimension-one stability up to near the Cauchy horizon.
\cite{DonningerOstermann2023}: global stability of the Yang--Mills profile,
odd $d$. \cite{Glogic2023}: global-in-space stability for Yang--Mills in
higher dimensions (preprint). \cite{Glogic2025}:
globally stable wave-map profile in all $d\ge3$. \cite{DonningerOstermann2024}:
wave maps and Yang--Mills in regions approaching the future light cone of
the singularity (preprint).
\item \emph{Optimal regularity and recent work.} \cite{DonningerWallauch2023}:
wave maps $\mathbb{R}^{1+4}\to S^4$, small data in the critical Sobolev
space. \cite{DonningerWallauch2025}: three-dimensional wave maps, fractional
critical space. \cite{DonningerMoscatelli2026a}: mode stability of
self-similar wave maps without symmetry for $d\ge4$ (preprint).
\cite{Gough2025,LiuRaees2025,GhoulLiuMasmoudi2025}: one-dimensional
derivative nonlinearities (Section~\ref{lit:acoustics}).
\item \emph{Complementary approaches.} Merle--Zaag: blowup rate
\cite{MerleZaag2003}; existence and universality of the profile in one
dimension \cite{MerleZaag2007}; in higher dimensions, the stability of
non-characteristic points and of the profile, and the dynamics near the
explicit (Lorentz-transformed) stationary solutions in similarity variables
\cite{MerleZaag2015,MerleZaag2016}. Bizo\'n et al.: numerics and heuristics
\cite{BizonChmajTabor2004}, countable families of self-similar profiles
\cite{BizonMaisonWasserman2007}. Lyapunov functionals for perturbed and
Klein--Gordon equations \cite{HamzaZaag2010,HamzaZaag2012,HamzaZaag2013},
and rates for NLKG \cite{KillipStovallVisan2014}. Speck: quasilinear
derivative-quadratic equations \cite{Speck2020}. Blowup surfaces via
Fuchsian/Nash--Moser methods \cite{Kichenassamy1996}, and scattering
constructions of stable ODE-type blowup on spacelike hypersurfaces
\cite{KadarLi2026} (preprint).
\end{itemize}

\subsection{Novelty assessment}
\label{lit:novelty}
We ran several targeted web searches per question (arXiv listings, journal
and repository pages, and citing papers), current to 13 October 2026.
Absence of evidence is not proof of absence; in particular the full texts
of \cite{KavallarisEtAl2015,Guo2010,Kichenassamy1996} were not accessible to
us.
\begin{enumerate}[leftmargin=1.4em,itemsep=2pt]
\item \emph{Hyperbolic MEMS $w_{tt}-\Delta w=-\lambda w^{-q}$ (damped or
not).} Existence, touchdown criteria and pull-in thresholds are rigorous
\cite{ChangLevine1981,FloresEtAl2007,Guo2010,KavallarisEtAl2011,Flores2014,KavallarisEtAl2015,Miyasita2017}.
The self-similar structure (ODE profile, non-trivial profiles, stability up
to time translation) is established only formally and numerically
\cite{GimperleinHeLacey2022,He2021}. We found no rigorous proof of the
quenching rate, of the profile, or of the stability of
\eqref{lit:eq:memsode}, in similarity coordinates or otherwise, for any
$q>1$ and any $d$. \textbf{Open}; a Donninger-type result would be new.
\item \emph{Exponential equation $u_{tt}-\Delta u=e^u$.} For smooth data
near constants, Kichenassamy's Nash--Moser construction \cite{Kichenassamy1996}
gives blowup on spacelike hypersurfaces in any dimension. In one dimension,
blowup curves and rates are known \cite{Godin2000,AzaiezMasmoudiZaag2016},
and the Liouville equation is explicitly solvable. We found no
Donninger-type theorem (stability of
$\log2-2\log(T-t)$ in $H^k\times H^{k-1}$ of the backward light cone with
exponential decay in $\tau$, $d\ge2$) and no Merle--Zaag-type higher-dimensional
profile result. \textbf{Partially known; the light-cone $H^k$ statement
appears open.}
\item \emph{Focusing Klein--Gordon with mass.} Blowup rates are known
\cite{KillipStovallVisan2014,HamzaZaag2013}. Stability of ODE blowup is
known for $m=0$ \cite{Ostermann2024}, and for strongly perturbed equations
in $d=3$ \cite{DonningerWallauch2021} (from the abstract we could not
confirm whether a mass term is included). Since $m^2u$ contributes an
$O(e^{-2\tau})$ term, the extension is expected to be routine, but we found
no explicit statement for general $d$ and $p$. \textbf{Expected, not found
explicitly.}
\item \emph{Damped power equation $u_{tt}+u_t-\Delta u=|u|^{p-1}u$.}
Blowup rates are known via Lyapunov functionals \cite{HamzaZaag2010,HamzaZaag2012},
and the one-dimensional profile in \cite{HamzaSaidi2021}. We found no
similarity-coordinate stability result for $d\ge2$. The damping term is an
$O(e^{-\tau})$ perturbation that is non-autonomous and breaks Lorentz
invariance. \textbf{Open in the $H^k$ light-cone sense, expected
perturbative.}
\item \emph{$u_{tt}-\Delta u=(\partial_tu)^2$.} In one dimension, stability
of the (generalised) self-similar blowup is proved
\cite{Gough2025,LiuRaees2025}. In three dimensions, stable ODE-type blowup
for a quasilinear class is proved by a different method \cite{Speck2020}.
A 2026 preprint (arXiv:2606.22282) treats radial blowup on spheres in
higher dimensions for a quadratic-derivative nonlinearity; we could not
verify its author list. % TODO verify authors of arXiv:2606.22282 before citing
\textbf{Largely known; a localised similarity-coordinate result for
$d\ge2$ was not found.}
\end{enumerate}
For the models of Sections~\ref{lit:heat} (full Cattaneo system) and
\ref{lit:acoustics} (Westervelt degeneracy) we found no singularity analysis
at all; both are quasilinear in the relevant regime and outside the scope of
the semilinear method.
